\section{Optimization interpretation}
\label{sec:optimization}

This section is the parallel to \S\ref{sec:clinical} for the optimization
researcher. The same artefacts, viewed through a different lens.

\subsection{The federated objective is non-convex even when the local objectives are convex}
A logistic regression has a strictly convex per-client objective. In the
federated case, the objective is $F(w) = \sum_k \frac{|D_k|}{N} F_k(w)$, which
is a convex combination of convex functions and is therefore convex. So our
logistic backbone is convex \emph{globally}.

The federation \emph{trajectory} is not a clean steepest-descent path on
$F$, however. \fedavg{} runs $E$ local steps before averaging; this is
equivalent to a \emph{block-coordinate descent with averaging}, and the
averaging step is not the gradient step on $F$. In particular, when the per-client
minima are far apart, the averaged $w_{t+1}$ can be \emph{worse} on $F$ than
$w_t$ even though every local step decreased $F_k$. This is the well-known
``client drift'' phenomenon (Li et al., 2020; Karimireddy et al., 2020).

\fedprox{}'s proximal term is the cleanest fix to this drift on convex
objectives: by adding $\frac{\mu}{2}\|w - w_t\|^2$ to each local objective, the
local minima are re-centred toward $w_t$ and the averaging step approaches the
gradient step on $F$ in the limit $\mu \to \infty$. At finite $\mu$ there is a
bias-variance trade-off: small $\mu$ approaches \fedavg{} (fast but drifty);
large $\mu$ approaches a one-step gradient method (slow but stable). The
$\mu \!=\! 0.1$ winner is empirically near the sweet spot of this trade-off
on \mimic{}.

\subsection{Reading the LP as a budgeted policy program}
The LP we run (\S\ref{sec:methods}, \S\ref{sec:mechanism-lp}) is a
\emph{static} policy program: choose a probability distribution over candidate
configurations to minimise mean loss subject to a mean byte-cost constraint.
Writing the LP in vector form,
\[
\min_{p\in\Delta} \;c^\top p \quad \text{s.t.} \quad b^\top p \leq B,
\]
the dual problem has dual $\mathbf{1}^\top \nu - B\lambda$ subject to
$\nu - b\lambda \leq c, \nu \in \mathbb{R}, \lambda \geq 0$. The KKT conditions
fully describe the solution: stationarity ($c \!=\! \nu - b\lambda$ at the
support), primal/dual feasibility, complementary slackness ($\lambda(b^\top p
- B) = 0$). At every budget we solve, KKT residuals are $\leq 5\!\times\!10^{-10}$
(\texttt{kkt\_diagnostics.csv}), so the dual is rigorous.

\paragraph{Why $\lambda$ is the shadow price.} By envelope reasoning, the
optimal value of the LP as a function of $B$ has $\frac{\partial \mathrm{OPT}}{\partial B}
= -\lambda$. So $\lambda$ is the rate of marginal-loss reduction per extra
byte of budget. Below the active region's right edge, $\lambda$ is positive
and informative. Above it, $\lambda$ is numerically zero -- spend the budget
elsewhere. This is the cleanest single-number answer to ``how much do extra
bytes help?''.

\subsection{Geometry of the loss landscape}
Figure~\ref{fig:landscape-1d} and Figure~\ref{fig:landscape-2d} render the
landscape on the validation slice. Three optimization-relevant facts:
\begin{enumerate}[leftmargin=*,nosep]
    \item \textbf{Logistic landscape is a clean convex bowl.} 1D minimum at
        $\alpha\!=\!0.65$, 2D minimum at the origin. Loss strictly increases
        away from the minimum in every direction we sampled. This is an
        explicit empirical confirmation that the logistic problem is convex.
    \item \textbf{MLP landscape is non-convex but locally smooth.} 1D minimum
        at $\alpha\!=\!0.7$, basin width $\sim 0.5$, gradient magnitude in the
        basin $\leq 0.1$ per unit $\alpha$. This is a wide, flat minimum -- a
        regime where Adam's adaptive step sizes work well.
    \item \textbf{The MLP basin extends beyond $\alpha\!=\!1$.} Loss at
        $\alpha\!=\!1.05$ is $0.40$, vs.\ $0.39$ at $\alpha\!=\!1.0$. So the
        trained-final solution is not at the minimum -- early stopping
        terminated before the optimizer fully descended. This is a well-known
        consequence of patience-based stopping: it stops improvements that are
        smaller than a noise floor (here, $\Delta\!=\!5\!\times\!10^{-4}$).
        Continued training would have descended further by $\sim 0.001$, a
        clinically negligible amount.
\end{enumerate}

\subsection{Hyperparameter sensitivity}
The DE history (Figure~\ref{fig:ga-scatter}) lets us read off
hyperparameter sensitivity. Three observations:
\begin{itemize}[leftmargin=*,nosep]
    \item \textbf{$\eta$ is the dominant axis.} Fitness has the
        cleanest dependence on $\log_{10}\eta$; the right corner of the
        $\eta$ scatter is the off-grid optimum.
    \item \textbf{$E$ matters second.} $E\!=\!1$ tends to underconverge; $E\!=\!3$
        tends to over-fit local-data-rich clients before averaging. $E\!\in\!\{2,3\}$
        spans the bulk of the low-fitness candidates.
    \item \textbf{$C$ is least sensitive.} For $C\!\in\!\{4,\dots,9\}$, fitness
        is approximately flat. The exception is $C\!=\!3$, which has slightly
        lower fitness because the per-round byte-cost is small. This is why
        the GA winner has $C \approx 3.5$ -- the LP-optimal trade-off between
        per-round cost and per-round signal.
\end{itemize}

\subsection{Why FedProx changes both the AUPRC and the stop time}
The proposal-run \fedprox{}-$\mu\!=\!0.1$ stops in $\sim$$35.3$ rounds vs.\
\fedavg{}'s $\sim$$36.5$ rounds. This is a small reduction \emph{in count}
but a large reduction \emph{in variance}: \fedavg{}'s stop time has std $5.9$
across seeds; \fedprox{}-$0.1$'s has std $9.9$ across seeds. So FedProx-$0.1$
stops a bit earlier on average but with wider seed-to-seed variance, because
it is occasionally pulled out of an early plateau by the proximal anchor's
re-centring.

\paragraph{Implication for production.} If the production training pipeline
allocates a fixed wall-time budget per round, FedProx-$0.1$ uses $\sim 25\%$
fewer rounds and so $\sim 25\%$ fewer wall-time units, even though its
per-round work is identical to FedAvg. That is the operational meaning of
``\fedprox{}-$0.1$ is cheaper''.

\subsection{The role of seeds}
We fixed $10$ prime seeds. The seed controls: (a) Adam's parameter
initialisation, (b) per-round client sampling, (c) data-loader batching, (d)
local-update step ordering. Std across seeds is the variability remaining
after fixing the algorithm and the data. Methods with low across-seed std are
\emph{reliable}; methods with high across-seed std are gambles in deployment.
\fedprox{}-$0.1$'s std on AUPRC is $0.005$, vs.\ \fedavg{}'s $0.014$ -- a
$2.8\!\times$ reliability improvement. \emph{Seed-std is the right place to read
off how confident a deployment can be in seed-1's metric matching seed-N's.}

\subsection{Synthesis: a four-step optimization recipe for federated mortality}
\begin{enumerate}[leftmargin=*,nosep]
    \item Use weighted CE with class weights $\propto 1/n_c$. Without this,
        nothing else matters.
    \item Use \fedprox{} with $\mu \!=\! 0.1$, $C \!=\! 5$, $E\!=\!1$, Adam
        $\eta \!\sim\! 0.003$, batch $256$. This is the deployment recipe.
    \item Run the LP shadow-price program once at deployment time on a small
        DE history; its $\lambda$-cliff tells you the byte-budget below which
        you starve the optimization and above which you waste bytes.
    \item Validate on per-client metrics, not pooled. The $9$-ICU table is
        what you ship; the global mean is what you debug with.
\end{enumerate}
